\section{Discussion}\label{sec:discussion}

The rational windows provide a computable map from programs to promised regimes. The conclusions do not require a decision criterion at noise rates or temperatures outside the windows. In the dynamical problem, the observed periodic signal gives a cycle of row laws for the time-one kernel; its phase average is stationary. The witness supplies nonconvergence, and the number of stationary row laws on the nonhalting side remains undetermined by this argument. The corresponding Gibbs states of the space--time interaction distinguish the phases separately because the DLR specification permits time-translation symmetry breaking. Blocking by the clock period turns them into translation-invariant states, and doubling adds an exact symmetry that exchanges them.

Several restrictions on the models remain natural to investigate. One is how small the alphabet and range can be for a single rule whose convergence diagram encodes halting. Another is which arrangements of halting and nonhalting windows can be obtained with one fixed automaton away from the small-noise regime. In the square-lattice problem, the interaction here is finite range and is built from an anchored pattern of one child and several parents. It would be useful to determine whether the same all-DLR dichotomy can be realized with only pair interactions at nearest-neighbor range and finite energies.


Theorems~D and~E characterize relatively effectively open sets as increasing unions of actual EWM sets of uniformly computable finite rational magnets. They leave a single-magnet realization problem: for rational \(0<a<b\), does every relatively effectively open \(W\subset[a,b]\) equal \(\mathcal M_\Phi\) for one finite rational magnet, with opposite translation-invariant ordered states on its complement? A sharp case asks for EWM on \([a,\beta_c)\) and order on \([\beta_c,b]\), where \(\beta_c=a+(b-a)\Omega\) and \(\Omega\in(0,1)\) is the halting probability of a fixed universal prefix-free machine. The lower cut of \(\Omega\) is effectively open. Purcell, Li and Cubitt constructed a family of two-dimensional translation-invariant nearest-neighbor quantum Hamiltonians with a single ground-state phase transition at a halting probability \cite{plc}; there the parameter enters the couplings, not the inverse temperature of one fixed classical magnet. Theorem~B already gives one magnet with a halting-programmed diagram on disjoint windows, and asserts nothing between those windows; it does not determine this cut.

Theorem~F settles the case of diagrams given explicitly as rational semialgebraic sets in \(z=e^{-\beta/d}\). The finite rational line-transfer mechanism used there has an EWM set semialgebraic in a Boltzmann coordinate. Its transition values are therefore exp-algebraic, meaning that \(e^{-\beta/d'}\) is algebraic for some positive integer \(d'\), and in particular they are computable. The algebraic temperature substitutions sketched in Section~\ref{sec:ceiling} would retain this limitation. Consequently this mechanism cannot place a transition at the noncomputable value \(a+(b-a)\Omega\). This is a ceiling of these finite-table constructions. It places no restriction on general finite magnets or on the programmed windows of Theorem~B. The alphabets and index lists produced by the synthesis may be extremely large; all searches terminate, and no elementary size bound is claimed.

\section*{How this work was done}
\addcontentsline{toc}{section}{How this work was done}
The assistance named on the first page was used in developing and checking the mathematical arguments. The proofs distinguish deterministic implications, finite probability estimates and explicitly stated external theorems. Independent mathematical checks informed the statements; exact finite calculations provide supplementary checks of the displayed examples. The authors take responsibility for the proofs, their scope and the use of the cited results.
